\section{Research Challenges and Opportunities}
\label{sec:challenges}
The comparisons in Sections~\ref{sec:tasks}--\ref{sec:evaluation} identify recurring questions about the decisions scientists retain, the transformations they must inspect, and the evidence supporting an interpretation. Feedback-driven classification raises questions about correction and propagated changes; generated operations raise questions about intent and execution provenance; shared and spatial workflows raise questions about authority, reference, and continuity. The following agenda develops these differences into design opportunities and testable questions concerning revisable delegation, inspectable transformations, traceable evidence, collaborative common ground, and continuity across environments.

The expanded cases make transformation risks concrete. Space-Time Hypercube projection can obscure spatial information and create apparent shrinkage as objects move away from a cutting plane; source-geometry views provide a partial safeguard.~\cite{ExpandedE22} PathwayMatrix's binary decomposition supports overview while limiting temporal and causal tracing.~\cite{ExpandedE37} Such trade-offs motivate preservation of source links and transformation assumptions, not rejection of computational abstraction.

\subsection{Delegation and Analyst Control}
\label{sec:challenges-delegation}
Delegated work must remain revisable when a scientist changes a question or discovers a problem. Approving an initial plan does not necessarily provide control over its later consequences. ProactiveVA illustrates this distinction: adjustable guidance frequency did not resolve a participant's request to intervene while an agent was executing.~\cite{EvalZhao2025} As Section~\ref{sec:evaluation-control} explains, control over assistance appearing and control over its execution are different capabilities.

A useful direction is to represent delegated work through explicit inputs, constraints, dependencies, and intervention points. If a scientist changes a cohort after inspecting a heatmap, the system should identify affected comparisons while preserving valid intermediate results. Routine recomputation under established parameters need not require the same supervision as changing the population or ranking criterion. Adaptive systems must additionally distinguish an inferred preference from permission to act; uncertainty and reversibility should influence whether to execute, clarify, or refrain.

Institutional deployment adds another level of control. REDCap's generative features are enabled by administrators through an institution-managed AI connection, with user permissions and researcher review of outputs; the underlying research data remain unchanged.~\cite{Shyr2026REDCapAI} This separates authority to enable a capability from responsibility for accepting its result. An analytical workflow should make both decisions visible and preserve links from an accepted summary to the records and request supporting it. The uptake evidence is discussed in Section~\ref{sec:evaluation-effectiveness}.

Studies should introduce changing goals during Selection, Filtering, and Precision Interaction and Comparison and Differentiation. Comparing whole-plan approval, operation-level approval, and consequential checkpoints would reveal whether analysts detect mismatches, redirect work, and preserve valid results. Total specification, supervision, and recovery effort matters more than the number of automated steps.

\subsection{Compression and Representational Fidelity}
\label{sec:challenges-compression}
Filtering, aggregation, and approximation make large datasets tractable while changing which relationships remain visible. The edge-tracing ambiguity discussed in Section~\ref{sec:evaluation-fidelity} illustrates why a clearer overview does not guarantee reliable detailed interpretation.~\cite{EvalHolten2006} Additional display space cannot recover evidence already removed from a representation.

A promising design direction connects summaries to their contributing observations and transformation choices, while indicating where reduction may matter. A heatmap of cell-group averages could expose group size and within-group variation so that a similar average does not silently conceal divergent responses. Access to detail is necessary but may be insufficient if the scientist has no reason to inspect the affected group. Adaptive aggregation also needs visible scope: a display change can reflect a changed representation rather than changed biology.

For Navigation and Multiscale Orientation and Comparison and Differentiation, evaluations should combine broad patterns with localized exceptions. Compare fixed summaries, expandable summaries, and summaries that indicate potentially obscured variation. Measure detection of exceptions, attribution of representational effects, and revision of conclusions when detail contradicts the overview. Compression succeeds scientifically when analysts can discover when omitted detail matters.

\subsection{Guidance and Diversity of Exploration}
\label{sec:challenges-guidance}
Predicting the analyst's next action is different from improving the inquiry, as the user-model comparison in Section~\ref{sec:evaluation-control} makes clear.~\cite{EvalHa2023} Guidance can reinforce familiar variables and explanations; the resulting interactions can then become evidence for further recommendations. Apparent preference stability may partly reflect the system's own influence.

Research should distinguish suggestions that extend an analysis, test its robustness, or examine competing explanations. In an expression study, recommending more genes with similar profiles can deepen one pattern without distinguishing treatment response from cell-composition differences. Matched-population comparisons or quality checks may offer more consequential alternatives. These purposes are interaction choices within existing assistance modes, not new taxonomy categories.

Systems could record whether feedback was volunteered or elicited by a suggestion and allow analysts to revise or suspend inferred preferences. Evaluations should assess recognition of conflicting evidence and justification of conclusions alongside coverage and effort. Focused exploration is not inherently biased, and more diverse exploration is not automatically better: the relevant question is whether assistance helps scientists recognize when their evidence warrants another direction.

\subsection{Generative Capability and Scientific Provenance}
\label{sec:challenges-generative}
Generated queries, code, views, and explanations introduce linked interpretive choices. Section~\ref{sec:evaluation-provenance} distinguishes a fluent account from an execution record; the PhenoFlow example also shows that correct execution cannot supply information absent from the data.~\cite{EvalKim2025} A revised question should retain the missing evidence or assumption that motivated the revision.

The research opportunity is to connect requests, plans, executed operations, results, and interpretations while recording their evidential status. CausalChat distinguishes suggested links from data-confirmed relationships; planning systems such as LightVA introduce dependencies through which an upstream revision can undermine later reasoning even if code still executes.~\cite{CausalChat,LightVA} These contextual examples motivate requirements rather than establish solutions for life-science practice.

A generated biological explanation should distinguish an observed difference, the comparison producing it, supporting literature, and the proposed mechanism. Changes to the comparison should identify dependent statements requiring reassessment. Studies can introduce unavailable variables, misleading explanations, execution discrepancies, and stale claims to test whether scientists locate the problem and correct it. Spoken or spatially presented assistance needs the same persistent route to evidence as a desktop explanation.

\subsection{Personalization and Collaborative Common Ground}
\label{sec:challenges-personalization}
Personalization can support different expertise while concealing differences in assumptions. The awareness--distraction tension in \emph{Talk to the Wall} motivates a related question: which consequential changes should collaborators see, and when?~\cite{EvalMolinaLeon2025} Its speech findings do not themselves validate personalized assistance; Section~\ref{sec:evaluation-collaboration} preserves that boundary.

Consider a biologist receiving population suggestions while a computational analyst receives normalization guidance. Private exploration may be useful, but changing a shared heatmap's underlying observations requires common awareness of the filter, its consequences, and whether the change has been adopted. Broadcasting every suggestion is unnecessary; silently changing shared evidence is unsafe for interpretation.

A useful direction supports transitions from private previews to attributed proposals and reversible shared changes. Permission to display a recommendation differs from permission to retrain a common model or publish a joint interpretation. Acceptance by the most active speaker should not imply collective agreement. Studies spanning separate and joint work should measure recognition of divergent assumptions, conflict resolution, participation, interruption costs, and the team's ability to explain a shared conclusion.

\subsection{Modality Specialization and Workflow Continuity}
\label{sec:challenges-continuity}
\label{sec:modality-continuity}
Moving between desktop, immersive, and shared environments should preserve object identities, selection criteria, transformations, uncertainty, and decision history while adapting controls to each setting. These workflows combine existing modalities. Navigation requires recovering context; selection preserves the intended subset; comparison preserves its basis; sensemaking traces evidence; collaboration makes consequential changes intelligible to others.

DashSpace provides an architectural example of cross-device synchronization while leaving semantic conflicts open.~\cite{AgendaBorowski2025} A tissue region selected in VR must refer to the same observations in a desktop matrix; a request about ``these cells'' needs a recorded referent. Spatial-language requirements and implemented scene operations provide complementary starting points.~\cite{SongEmbodiedNLI,VOICE} Research should expose ambiguity, invalidate stale references, and make changes to shared analysis attributable and reversible.

Preferences inferred in one setting should not silently govern a different task. Accessible alternatives should permit inspection and correction without a headset or particular input method. Studies should follow complete transitions, examining reorientation, stale-state detection, correction, and participation, and tracking who proposes, scopes, performs, and assesses consequential operations.

\subsection{Shared Scenarios and Analytical Traces}
\label{sec:challenges-evaluation}
Shared life-science scenarios and analytical traces could support the evaluations developed in Section~\ref{sec:evaluation}. Domain experts should document plausible questions, alternatives, errors, and recovery paths, including inappropriate filters, contradictory evidence, and changing goals. Traces should connect requests, data versions, parameters, transformations, outputs, and revisions, recording who proposed, scoped, performed, and assessed each operation and the opportunities to inspect, change, or undo it. Shared definitions, trace schemas, and permitted data or synthetic substitutes would allow assistance ablations and matched workflow studies to reuse comparable conditions while accommodating several defensible conclusions in open-ended inquiry.
